\documentclass[11pt]{article}
\usepackage[T1]{fontenc}
\usepackage[utf8]{inputenc}
\usepackage{lmodern}
\usepackage[margin=0.86in]{geometry}
\usepackage{amsmath,amssymb,amsthm,booktabs,longtable,array}
\usepackage[hidelinks]{hyperref}
\usepackage{enumitem}
\usepackage{microtype}
\usepackage{url}

\newtheorem{theorem}{Theorem}
\newtheorem{proposition}{Proposition}
\newcommand{\R}{\mathbb R}
\newcommand{\RobInv}{\mathsf{RobInv}_{\mathrm{src}}}
\newcommand{\EqSrc}{\mathsf{EqSrc}_{\mathrm{prot}}}
\newcommand{\status}[1]{\textsf{#1}}
\newcolumntype{P}[1]{>{\raggedright\arraybackslash}p{#1}}
\setlist{nosep,leftmargin=*}
\setlength{\emergencystretch}{2em}
\urlstyle{tt}

\title{Refuter Stress of Source-Law-Space\\
Robust-Invariance Protection}
\author{Refuter parent synthesis, RT-20260811-009}
\date{11 August 2026}

\begin{document}
\maketitle

\begin{center}
\fbox{\begin{minipage}{0.94\linewidth}
\textbf{Decisive result:} \status{scoped\_obstruction}.
The fixed RT007 theorem is correct on its declared proposal-only roots, and
the RT008 \status{source\_pure\_as\_written} verdict remains correct at its
narrow written-syntax scope.  A fixed positive-margin tuple has an open
robustness neighborhood.  What fails is the stronger use of the unchanged
predicate as a source-derived selector and family-wide P4--T02 protection
bridge: neither the predicate nor current ontology selects the root family or
protects a uniform margin, variation bound, full-state lift, orientation, and
nontrivial presentation groupoid across it.

\smallskip
\textbf{Route consequence:} replay of that unchanged standalone use is
locally frozen.  Conditional theorem, certificate, and diagnostic use remain
open.  One Theoretical Continuation Selector packet is selected but not
executed.  Six local freezes are active; every Distance-to-GR row remains
\status{no\_delta}; D7 is not reevaluated, B2 is inactive, and P4--T03 is
locked.
\end{minipage}}
\end{center}

\section{Fixed candidate and exact scope}

RT007 fixes the proposal-only datum
\[
 \mathcal D_{\mathrm{prot}}
 =(X_{\mathrm{law}},K,\mathcal A,\Delta,\EqSrc)
\]
and defines strict robust invariance by requiring every admitted velocity to
lie in the interior of the contingent tangent cone on each relevant boundary
stratum.  Under the stated regularity and locally uniform positive-margin
hypotheses, RT007 proves strong forward invariance, a finite conormal
certificate, positive certificate regraduation, coherent complete-tuple
presentation transport, and a balanced-normal obstruction.  RT008 finds no
hidden target import in those written premises while leaving every root
independently unprovenanced and unadopted.

This Refuter packet holds those results fixed.  It attacks only the overread
that the forward predicate selects or protects its own complete family of
roots.  A failure of that overread is not a counterexample to the conditional
theorem.

\section{Positive control: a fixed margin is locally stable}

Let the active residual of a fixed regular datum be
\[
 r_{i,F,E}(x)=d\rho_i(x)[F(x)+E].
\]

\begin{theorem}[Fixed-margin perturbation control]
Suppose every relevant active residual is bounded below by one number
\(m>0\).  If an admitted field perturbation \(W\) obeys
\[
 \left|d\rho_i(x)[W(x)]\right|\leq\delta<m
\]
uniformly over the same active tests, then every perturbed residual is at
least \(m-\delta>0\).  Hence the same fixed datum remains strictly robustly
invariant under the RT007 hypotheses.
\end{theorem}

\begin{proof}
For each active test,
\[
 d\rho_i[F+E+W]
 \geq d\rho_i[F+E]-|d\rho_i[W]|
 \geq m-\delta>0.
\]
The RT007 active-conormal characterization and strong-invariance theorem then
apply without modification.
\end{proof}

This positive control is essential.  The obstruction below does not claim
that an arbitrarily small perturbation destroys every fixed passing tuple.
The permitted radius is tied to the supplied margin, and the proposal does
not derive a positive lower bound for a selected root family.

\section{Margin collapse and variation enlargement}

On \(X=\R\), take \(K=\R_{\geq0}\), the inward conormal \(dx\),
\(\Delta=\{0\}\), and the constant generator
\[
 F_\epsilon(x)=\epsilon,
 \qquad \epsilon\geq0.
\]

\begin{proposition}[Open fixed members, unprotected family boundary]
Every \(\epsilon>0\) gives a strict pass with exact margin \(\epsilon\),
whereas \(\epsilon=0\) fails strict inclusion.  Moreover
\(F_\epsilon\to F_0\) uniformly in every compact-open \(C^k\) seminorm.
Thus each passing member is locally open, but the family has no positive
uniform protection radius.
\end{proposition}

\begin{proof}
At the only boundary point, \(dx[F_\epsilon]=\epsilon\).  This is positive
exactly when \(\epsilon>0\).  All positive-order derivatives of the constant
fields agree, and the zeroth-order difference is exactly \(\epsilon\) on
every set.
\end{proof}

Now use the nested compact variation family
\[
 \Delta_\alpha=[-\alpha\epsilon,+\alpha\epsilon],
 \qquad \alpha\geq0,\quad \epsilon>0.
\]
Its worst residual is
\[
 \inf_{E\in\Delta_\alpha} dx[F_\epsilon+E]
 =(1-\alpha)\epsilon.
\]
Therefore \(\alpha<1\) strictly passes, \(\alpha=1\) reaches zero-margin
strict failure, and \(\alpha>1\) contains an outward velocity.  The
classifier faithfully detects the threshold, but it does not select the
admissible value of \(\alpha\) or prove that current ontology bounds the
variation class below it.

\section{Orientation transport is not orientation selection}

The original tuple \(K=\{x\geq0\}\), \(F=+1\), with inward conormal
\(+dx\), has residual \(+1\).  Under the reflection \(y=-x\), complete
tuple transport gives
\[
 K'=\{y\leq0\},\qquad F'=-1,\qquad \text{inward conormal }-dy,
\]
and \((-dy)(-1)=+1\).  The pass is preserved exactly.

By contrast, replacing the admissible set with \(\{y\leq0\}\) while holding
the generator fixed at \(+1\) gives \((-dy)(+1)=-1\).  That operation is not
complete presentation transport; it is a different root tuple.  Covariance
therefore preserves a supplied orientation but cannot choose it.

\section{Same-slice full-state lifts have opposite verdicts}

Let
\[
 K'=\R_{\geq0}\times[-\rho,\rho],
 \qquad \epsilon,\nu,\rho>0,
\]
and compare
\[
 F^{-}(x,z)=(\epsilon,-\nu z),
 \qquad
 F^{+}(x,z)=(\epsilon,+\nu z).
\]
Both agree with \((\epsilon,0)\) on the embedded slice \(z=0\).

\begin{theorem}[Full-state lift nonselection]
The lift \(F^{-}\) strictly protects all three boundary types of \(K'\),
whereas \(F^{+}\) fails at both transverse faces.  Hence the fixed reduced
slice and its robust-invariance verdict do not select a full-state lift.
\end{theorem}

\begin{proof}
At \(x=0\), both fields have residual \(\epsilon>0\).  At the lower face
\(z=-\rho\), the inward conormal is \(+dz\); at the upper face
\(z=+\rho\), it is \(-dz\).  The two transverse residuals of \(F^{-}\)
are \((+\nu\rho,+\nu\rho)\), while those of \(F^{+}\) are
\((-\nu\rho,-\nu\rho)\).
\end{proof}

The passing lift has its own local protection threshold.  For
\[
 F_{\lambda,\eta}(x,z)=(\epsilon,-\lambda z+\eta)
\]
the lower and upper transverse residuals are
\(\lambda\rho+\eta\) and \(\lambda\rho-\eta\).  Strict protection holds
exactly when \(|\eta|<\lambda\rho\), reaches zero at equality, and fails
beyond it.  With \(\lambda=0\), the two residuals are \(\eta\) and
\(-\eta\), so no constant transverse drift strictly protects both faces.
Dimensional extension therefore requires new supplied inward-law structure;
the predicate itself does not construct it.

\section{Nontrivial presentation arrows are extra structure}

Consider two source objects, each with \(X=\R\) and \(K=\R_{\geq0}\).
The first has the singleton generator \(F_0=1\); the second has
\(F_1=2\).  Both objects strictly pass.  Their identity-only groupoids are
automatic.

An identity-coordinate map between the two objects is not a presentation
arrow because its derivative sends \(F_0\) to \(1\), not \(2\).  A scaling
map \(g(x)=2x\), however, does send \(F_0\) to \(F_1\), and its inverse sends
\(F_1\) back to \(F_0\).  Thus classifier truth determines neither the
existence nor the intended class of nontrivial arrows.  Coherence is a valid
conditional theorem once the groupoid is supplied, not a source-side selector
of the groupoid.

The RT008 dimension-one and dimension-two passing tuples supply a further
control: a common pass value cannot place them in one \(C^1\) presentation
orbit because ambient dimension is invariant.

\section{Decisive obstruction and freeze}

The exact result is a candidate-scoped obstruction:
\begin{quote}
Neither \texttt{SourceLawSpaceRobustInvarianceProtection\_v1} nor current
ontology selects a complete root family and proves a uniform positive
conormal margin, independently admissible variation bound, transverse closure
law, orientation transport, and nontrivial \(\EqSrc\) coherence across that
family.
\end{quote}

The obstruction identifier is
\begin{flushleft}
\small\ttfamily
OB-V22-P4T02-B2-SOURCE-LAW-SPACE-ROBUST-INVARIANCE-\\
PROTECTION-ROBUSTNESS-001.
\end{flushleft}
It freezes only unchanged predicate replay as a standalone source-derived
selector or family-wide P4--T02 bridge, under
\begin{flushleft}
\small\ttfamily
NDCL-V22-P4T02-B2-SOURCE-LAW-SPACE-ROBUST-INVARIANCE-\\
PROTECTION-ROBUSTNESS.
\end{flushleft}
Conditional theorem, certificate, exact fixtures, and diagnostic use remain
available.  A materially distinct source-side theorem, witness, bridge, or
protected human-gated ontology decision remains possible.

The five inherited freezes also remain exact:
\begin{itemize}
  \item \path{NDCL-V22-P4T02-B2-SHARED-TAU-SELECTOR};
  \item \path{NDCL-V22-P4T02-B2-REPAIRED-QUOTIENT-DESCENT-ROBUSTNESS};
  \item \path{NDCL-V22-P4T02-B2-COMMON-CHARACTER-HOLONOMY-PROTECTION};
  \item \path{NDCL-V22-P4T02-B2-ORIENTED-MATROID-BRIDGE-SELECTION-ROBUSTNESS};
  \item \path{NDCL-V22-P4T02-B2-SIGNED-CUBIC-VIABILITY-SELECTOR-ROBUSTNESS}.
\end{itemize}

\section{Distance-to-GR}

\begin{longtable}{P{0.29\linewidth}P{0.61\linewidth}}
\toprule
Burden & RT009 status \\
\midrule
\endfirsthead
\toprule
Burden & RT009 status \\
\midrule
\endhead
Source ontology primitives & No delta; no complete root family or family-protection law is selected or adopted. \\
Source equivalence \(\mathsf{EqSrc}\) & No delta; nontrivial arrows and coherence remain supplied conditional data. \\
\(\mathsf{RetainH}\) & No delta; no primitive or theorem. \\
\(\mathsf{GenH}\) & No delta; generator families and perturbation classes remain proposal-only inputs. \\
\(\mathsf{ObsLoc}_{lc}\) & No delta; no observer-localization or operational readout map. \\
\(\mathsf{Resp}_{lc}\) & No delta; invariance verdicts and certificates are not empirical responses. \\
\(M_{\rm src}\) & No delta; source-law charts and product controls are not adopted physical source spacetime. \\
\(g_{\rm eff}\) & No delta; no physical cone, conformal class, scale, or effective metric. \\
Matter coupling & No delta; no declared-sector universal propagation or coupling law. \\
Einstein equations & No delta; no action, variation, stress energy, or field equation. \\
Finite-variation robustness & No delta; fixed margins are locally stable, but family-uniform protection and full-state closure are not source-derived. \\
Benchmark promotion & No delta; no exact-GR comparison, recovery, or promotion. \\
Gate Chair status & No delta; no protected adoption or verdict. \\
Current route freeze or hard-fail & No delta; five inherited freezes remain active and unchanged robust-invariance protection replay is locally frozen. \\
\bottomrule
\end{longtable}

\section{Next bounded packet}

Exactly one future packet is selected but not executed:
\begin{flushleft}
\small\ttfamily
PKT-V22-P4T02-B2-POST-SOURCE-LAW-SPACE-ROBUST-INVARIANCE-\\
REFUTER-THEORETICAL-CONTINUATION-SELECTION-V1.
\end{flushleft}
Its type, next role, and status are, respectively,
\begin{flushleft}
\small\ttfamily
theoretical\_continuation\_selector,\\
theoretical-continuation-selector@0.1.0,\\
selected\_not\_executed.
\end{flushleft}
It must select one materially distinct
same-milestone packet without replaying any of the six frozen standalone
candidates.

\section{Explicit non-conclusions}

This task does not refute the RT007 fixed-domain theorem or reverse RT008
written-syntax purity.  It does not repair, adopt, or reject the candidate or
execute the successor.  It does not reevaluate D7 or P2--T01 adequacy, act on
B2, complete P4--T02 for plan dependency, unlock P4--T03, or construct a
physical causal cone, time orientation, empirical response, universal matter
propagation, conformal geometry, physical scale, \(g_{\rm eff}\), coupling,
or Einstein equations.  It creates no canonical ontology, Gate Chair,
benchmark, proof, publication, release, push, external-action, global-no-go,
future-extension-impossibility, or completed-derivation authority.

\end{document}
